\documentclass[a4paper, titlepage]{book}

\usepackage{amsmath}
\numberwithin{equation}{section}
\usepackage{amssymb}
\usepackage{csquotes}
\MakeOuterQuote{"}
\usepackage[useregional]{datetime2}
\DTMsetstyle{iso}
\usepackage{hyperref}
\hypersetup{
    colorlinks = true,
    allcolors = blue
}
\usepackage{indentfirst}
\usepackage[skip=5pt, indent=30pt]{parskip}
\usepackage{pgf}
\usepackage{physics2}
\usepackage{siunitx}
\sisetup{
    input-digits = 0123456789\pi,
    inter-unit-product = \cdot,
}


\title{
    Eric Zhang's Physics Notes \\
    \large Volume 2: Electricity and Magnetism
}
\author{Eric Zhang}
\date{\today}


\begin{document}
\frontmatter
\maketitle
\tableofcontents


\mainmatter

\chapter{The Electric Force}

\textbf{Electric charge} is denoted by the symbol \(q\) or sometimes \(Q\),
and its SI unit is the coulomb (C).

Charge can be either positive or negative,
with protons carrying a positive charge and electrons carrying a negative charge.

The \textbf{Law of Conservation of Charge} states that charge can not created or destroyed,
only transferred,
and the net charge has to remain constant.
There are two ways to charge and discharge (i.e., move charges):
\begin{itemize}
    \item Conduction -- friction (triboelectricity) or contact
    \item Induction -- contactless
\end{itemize}

Charged objects exert an \textbf{electric force} on other charged objects around them.
The magnitude of the electric force can be given by \textbf{Coulomb's Law}:
\begin{equation}
    F = k\frac{q_1 q_2}{r^2}
\end{equation}
where
\begin{itemize}
    \item \(F\) is the force \emph{both} objects experience as given by Newton's Third Law.
    \item \(k \approx \qty{8.99e-9}{\newton \meter\squared \per \coulomb\squared}\)
          and is called the Coulomb constant.
    \item \(q_1\) and \(q_2\) are the two charges of the two objects.
          Note that both values are signed.
    \item \(r\) is the distance between the two particles.
\end{itemize}
Objects with charges of the same sign will repel each other,
while objects with charges of opposite signs will attract each other.
Thus, Coulomb's Law will give a \emph{negative} value for attraction
and a \emph{positive} value for repulsion.

\chapter{Maxwell's Equations}


\appendix
\chapter{Mathematics and Notation}

\section{Vectors} \label{section:vectors}

\subsection{Vectors vs. Scalars}

\backmatter
\chapter{References}


\end{document}